\bigskip
\bigskip
\subsubsection*{Խնդիրներ և հարցեր թեմա 16-ի վերաբերյալ}

\begin{enumerate}[label=\thesection.\arabic*.]


% 16.1
\item Ճի՞շտ է պնդումը․ ցանկացած անտիդիսկրետ տարածություն գծորեն կապա\-կց\-ված է:

% 16.2
\item Ապացուցեք․ գծորեն կապակցված ամեն մի $X$ տարածության ցանկացած $\faktor{X}{R}$ ֆակտոր-տարածություն գծորեն կապակցված է:

% 16.3
\item Ճի՞շտ է պնդումը․ եթե տոպոլոգիական տարածությունների $X \times Y$ արտադրյալը գծորեն կապակցված է, ապա $X$-ը և $Y$-ը գծորեն կապակցված են:
\begin{hint}
Դիտարկեք $X \times Y$ արտադրյալի $P_X$ և $P_Y$ կանոնական պրո\-յեկ\-ցի\-ա\-ները $X$-ի և $Y$-ի վրա և օգտվեք թեորեմ 5-ի ապացույցից:
\end{hint}

% 16.4
\item Ապացուցեք հետևյալ թեորեմը․ $X \times Y$ արտադրյալը գծորեն կապակցված է այն և միայն այն դեպքում, երբ գծորեն կապակցված են $X$-ը և $Y$-ը:
\begin{hint}
Պայմանի բավարարությունը ապացուցելու համար դիտարկենք $X \times Y$-ի որևէ երկու $(x_1, y_1)$ և $(x_2, y_2)$ կետեր։ Դիցուք $f_1: I \to X$ ուղին մի\-աց\-նում է $x_1$ կետը $x_2$ կետին, իսկ $f_2: I \to Y$ ուղին միացնում է $y_1$ կետը $y_2$ կետին։ Օգտվելով թեմա 14-ում թեորեմ 3-ից՝ դիտարկեք $(f_1, f_2): I \to X \times Y$, $(f_1, f_2)(t) = \big(f_1(t), f_2(t)\big)$ արտապատկերումը:
\end{hint}

% 16.5
\item Դիցուք $X$-ը և $Y$-ը հոմեոմորֆ տարածություններ են։ Ապացուցեք․ $X$-ը գծո\-րեն կապակցված է այն և միայն այն դեպքում, երբ գծորեն կապակցված է $Y$-ը:

% 16.6
\item Ապացուցեք․ $\R^n$ էվկլիդեսյան տարածության $B^n = \left\{x;\ \|x\| \le 1\right\}$ միավոր գուն\-դը գծորեն կապակցված ենթատարածություն է:
\begin{hint}
Հիմնավորեք, որ $B^n$-ը ուռուցիկ ենթաբազմություն է:
\end{hint}

% 16.7
\item Ապացուցեք, որ $\R^{n+1}$ էվկլիդեսյան տարածության $S^n_+$ և $S^n_-$ միավոր կիսա\-սֆե\-րաները, որտեղ
\[ 
S^n_+ = \left\{x \in \R^{n+1};\ \|x\|=1,\, x_{n+1} \ge 0\right\}, \quad S^n_- = \left\{x \in \R^{n+1};\ \|x\|=1,\, x_{n+1} \le 0\right\}
\]
գծորեն կապակցված ենթատարածություններ են:
\begin{hint}
Հիմնավորելով, որ $\R^{n+1}$-ի պրոյեկցիան $\R^n$-ի վրա՝ սահմանված $P(x_1, x_2, \dots, x_{n+1}) = (x_1, x_2, \dots, x_n)$ բանաձևով, որոշում է $S^n_+$ և $S^n_-$ կիսա\-սֆե\-րաներից յուրաքանչյուրի հոմեոմորֆիզմ $B^n$ գնդին, օգտվեք նախորդ երկու խնդիր\-ներից:
\end{hint}

% 16.8
\item Ապացուցեք․ $\R^{n+1}$ էվկլիդեսյան տարածության $S^n = \left\{x;\ \|x\|=1\right\}$ միավոր սֆերան գծորեն կապակցված տարածություն է:
\begin{hint}
Ներկայացնելով $S^n$-ը որպես $S^n_+$ և $S^n_-$ կիսասֆերաների միավո\-րում՝ օգտվեք նախորդ խնդրից և թեորեմ 3-ից:
\end{hint}

% 16.9
\item Ապացուցեք․ $n \ge 0$ դեպքում $\mathbb{RP}^n$ իրական պրոյեկտիվ տարածությունը գծո\-րեն կապակցված է:
\begin{hint}
Երբ $n=0$, $\mathbb{RP}^0$-ն բաղկացած է մի կետից։ Երբ $n > 0$, ներկա\-յաց\-րեք $\mathbb{RP}^n$-ը որպես $S^n$ սֆերայի կամ $S^n_+$ կիսասֆերայի ֆակտոր-տարածու\-թյուն և օգտվեք նախորդ խնդիրներից:
\end{hint}

% 16.10
\item Դիտարկենք $\R^2$ հարթության 
\[
A = \left\{(x, y);\ x > 0,\, y = \sin \dfrac{1}{x}\right\} \quad  \text{ և } \quad B = \left\{(0, y);\ -1 \le y \le 1\right\}
\] 
ենթաբազմությունները։ Ապացուցեք․ $A \cup B$ միավորումը կապակցված է, բայց գծորեն կապակցված չէ:
\begin{hint}
Թեորեմ 7-ի նմանությամբ ցույց տվեք, որ $B$-ի ոչ մի կետ հնա\-րա\-վոր չէ ուղիով միացնել $A$-ի որևէ կետի հետ:
\end{hint}

\end{enumerate}
